\snsubsection{Convertidor reductor o Buck}

\begin{figure}[H]
    \centering
    \begin{tikzpicture}
        % Paths, nodes and wires:
        \draw (4.625, 15.125) to[american voltage source] (4.625, 13.125);
        \draw node[pigfete, solderdot, rotate=90] at (6.395, 15.625) {};
        \draw (7.875, 15.625) to[american inductor] (9.875, 15.625);
        \draw (7.625, 13.125) to[empty Schottky diode] (7.625, 15.125);
        \draw (10.125, 13.125) to[capacitor] (10.125, 15.125);
        \draw node[circ] at (7.625, 15.625) {};
        \draw node[circ] at (7.625, 15.625) {};
        \draw node[circ] at (7.625, 12.625) {};
        \draw node[circ] at (7.625, 12.625) {};
        \draw node[circ] at (10.125, 15.625) {};
        \draw node[circ] at (10.125, 15.625) {};
        \draw node[circ] at (10.125, 12.625) {};
        \draw node[circ] at (10.125, 12.625) {};
        \draw (7.625, 15.125) -| (7.625, 15.625);
        \draw (7.625, 13.125) -| (7.625, 12.625);
        \draw (10.125, 12.625) -| (10.125, 13.125);
        \draw (10.125, 15.125) -| (10.125, 15.625);
        \draw (4.625, 15.125) |- (5.125, 15.625) -- (5.625, 15.625);
        \draw (4.625, 13.125) -- (4.625, 12.625);
        \node[shape=rectangle, draw, line width=1pt, inner sep=0, minimum width=1.215cm, minimum height=0.965cm](Rect1) at (6.125, 13.625){} node[anchor=center] at (Rect1.center){$\mu C$};
        \draw node[circ] at (6.125, 14.125) {};
        \draw (6.125, 14.645) |- (6.125, 14.125);
        \node[shape=rectangle, draw, line width=1pt, inner sep=0, minimum width=1.215cm, minimum height=1.215cm](Rect2) at (14.375, 14){} node[anchor=center] at (Rect2.center){$Carga$};
        \draw node[circ] at (14.375, 14.625) {};
        \draw node[circ] at (14.375, 13.375) {};
        \draw (11.625, 15.625) -| (14.375, 14.625);
        \draw (11.625, 12.625) -| (14.375, 13.375);
        \draw (4.625, 12.625) -- (10.125, 12.625);
        \draw (7.165, 15.625) -- (7.875, 15.625);
        \draw (9.875, 15.625) -- (10.125, 15.625);
        \draw (11.625, 15.625) to[american resistor] (11.625, 14.125);
        \draw (11.625, 14.125) to[american resistor] (11.625, 12.625);
        \draw (11.625, 15.625) -- (10.125, 15.625);
        \draw (11.625, 12.625) |- (10.125, 12.625);
        \draw node[circ] at (11.625, 15.625) {};
        \draw node[circ] at (11.625, 12.625) {};
        \draw node[circ] at (11.625, 14.125) {};
        \node[shape=rectangle, draw, line width=1pt, inner sep=0, minimum width=1.215cm, minimum height=0.965cm](Rect3) at (12.875, 13.375){} node[anchor=center] at (Rect3.center){$Retro$};
        \draw (11.625, 14.125) -- (12.375, 14.125) -| (12.875, 13.875);
        \draw node[circ] at (12.875, 13.875) {};
    \end{tikzpicture}
    \caption{Circuito básico del convertidor Buck.}
    \label{fig:buck}
\end{figure}

Un convertidor Buck básico consiste en una conmutación controlada de un interruptor, un diodo, un filtro de salida, y un circuito de realimentación. El interruptor es el encargado de traspasar la potencia de entrada a la salida, dependiendo de la velocidad de conmutación o del tiempo que esté encendido, la tensión en la salida puede ser aumentar o disminuir, manteniéndose siempre por debajo de la tensión de entrada.

El voltaje medio de salida del convertidor puede ser controlado de 2 formas, PWM (pulse width modulation) o PFM (pulse frequency modulation). En el primero de los casos lo que se modifica es el tiempo de encendido, manteniéndose el periodo constante. A diferencia de PFM donde el periodo es el que se modifica (cambia la frecuencia de conmutación), pero el tiempo de encendido se mantiene constante, por ejemplo en un 50 % del duty. El método más elegido para los convertidores buck es el PWM.

Otro punto importante de los convertidores es el modo de operación, el mismo puede operar en dos modos de conducción, CCM (continuous conduction mode) y DCM (discontinuous conduction mode); en el primero la corriente en el inductor de salida permanece en un valor positivo, haciendo que nunca tenga corriente cero a través del mismo en el periodo de conmutación. En cambio para el caso de DCM, la corriente en el inductor se hace cero en algunos momentos del periodo de conmutación.

\begin{figure}[H]
    \centering
    \includegraphics[width=0.8\textwidth]{buck/CCM-DCM-Graphs.png}
    \caption{Modos de funcionamiento.}
    \label{fig:CCM-DCM}
\end{figure}

\snsubsubsection{Modo encendido}
\begin{figure}[H]
    \centering
    \includegraphics[width=0.8\textwidth]{buck/buck-encendido.png}
    \caption{.}
    \label{fig:buck-encendido}
\end{figure}

En este momento el diodo queda polarizado de forma inversa, haciendo que toda la corriente de entrada se aplique en el inductor; el inductor se carga.

\begin{figure}[H]
    \centering
    \includegraphics[width=0.5\textwidth]{buck/buck-waveform.png}
    \caption{.}
    \label{fig:buck-waveform}
\end{figure}

En la figura anterior se observa que, mientras el interruptor (en la mayoría de casos un transistor) conduce, la corriente a través del inductor aumenta llegando a su valor máximo; en el momento en que se apaga el interruptor la corriente comienza a disminuir.

\snsubsubsection{Modo apagado}
\begin{figure}[H]
    \centering
    \includegraphics[width=0.8\textwidth]{buck/buck-apagado.png}
    \caption{.}
    \label{fig:buck-apagado}
\end{figure}

En este momento, el interruptor se encuentra abierto, y el inductor cambia su polaridad para oponerse a la descarga que está sufriendo el mismo, y el diodo queda polarizado en forma directa.